\section{Proof Production and Minimization}
\label{sec:proofs}

% Source: ../egg-papers/2022-explanations-fmcad/01-intro.tex;
% Herbie proof documentation supplies the explanation application.
\sourcenote{Edited from \emph{Small Proofs from Congruence Closure}, \S1;
  explanation and debugging motivation added.}
Small proofs reduce both validation costs and the number of steps a person must inspect
  to understand or debug an optimizer's result~\cite{herbie-proofs}.
Our work, published at FMCAD 2022, retains alternative explanations during congruence closure
  and presents two algorithms for choosing among them~\cite{eggproof}.
\textsc{TreeOpt} is the theoretical contribution, optimizing a relaxed proof-size objective;
  greedy proof search is the practical algorithm implemented in \texttt{egg}.

\subsection{Background: Congruence Proof Generation}
\label{sec:background-proof-production}

% Source directory: ../egg-papers/2022-explanations-fmcad/.
% Sources: 02-background.tex and 03-explanation-length.tex.
\sourcenote{Edited from \emph{Small Proofs from Congruence Closure}, \S\S2--3.}
A \emph{proof graph} records justifications for equalities:
  its vertices are terms and its edges connect equal terms.
An edge can record an \emph{input equality}, a concrete equality supplied to congruence closure,
  for example by a rewrite.
Alternatively, a \emph{congruence edge} justifies $f(a)=f(b)$ by referring to a proof of $a=b$.
A proof follows a path between its endpoints, combines the edges by transitivity of equality,
  and recursively proves the child equalities required by congruence edges.
A rewrite rule may supply many different input equalities,
  and a proof may use the same input equality more than once.

% Sources: same files, spanning-tree construction and discarded equalities.
\sourcenote{Edited from \emph{Small Proofs from Congruence Closure}, \S\S2--3.}
Traditional proof generation maintains a tree connecting the terms in each e-class~\cite{pp-congr}.
Its unique paths give valid proofs, but additional equalities between already connected terms
  are discarded even when they offer shorter explanations.
Our approach retains those equalities as alternative edges
  without changing the asymptotic cost of congruence closure.

\subsection{Generating Small Proofs}
\label{sec:small-proofs}

% Sources: Small Proofs, 03-explanation-length.tex and 04-flattened-optimal.tex.
\sourcenote{Edited from \emph{Small Proofs from Congruence Closure}, \S\S3--4.}
Our \emph{DAG size} metric counts distinct input equalities,
  counting an equality once even if several subproofs use it.
The name refers to representing a proof as a directed acyclic graph (DAG) with shared subproofs.
Finding a proof of minimum DAG size is NP-hard~\cite{small-conflict}.
Sharing makes local choices difficult: the cost of a subproof depends on
  which equalities other parts of the proof already use.
Our relaxed \emph{proof tree size} metric counts every use of an input equality
  after expanding all congruence justifications.

% Source: Small Proofs, 02-background.tex and 04-flattened-optimal.tex;
% inference-tree layout and explanation are new for the report.
\begin{samepage}
\sourcenote{New inference-tree presentation of the example from
  \emph{Small Proofs from Congruence Closure}, \S\S2, 4.}
For example, $a+0=a$ and $2+2=4$ give this proof:
\[
  \dfrac{
    a+0=a
    \qquad
    \dfrac{a+0=a\qquad 2+2=4}{g(a+0,2+2)=g(a,4)}\;\text{congruence}
  }{f(a+0,g(a+0,2+2))=f(a,g(a,4))}\;\text{congruence}.
\]
\end{samepage}
Each horizontal line applies congruence to the equalities above it to derive the equality below it.
The proof uses $a+0=a$ twice and $2+2=4$ once, so its tree size is three and its DAG size is two.
Tree size therefore bounds DAG size from above.
Because it counts repeated uses separately, the cost of each subproof
  is independent of the equalities used elsewhere.
This lets us find a proof of minimum tree size in polynomial time.

\subsubsection{TreeOpt: Optimal Proof Tree Size}
\label{sec:treeopt}

% Source: ../egg-papers/2022-explanations-fmcad/04-flattened-optimal.tex,
% Minimum Proof Tree Size Algorithm; notation replaced with explanation.
\sourcenote{Edited from \emph{Small Proofs from Congruence Closure}, \S4.}
If every edge had a known cost, finding a small proof would be a shortest-path problem:
  choose a path whose edge costs have the smallest sum.
An input-equality edge costs one, but a congruence edge costs the tree size
  of the subproofs that justify it.
Those subproofs may themselves use congruence,
  so their costs are initially unknown.

% Source: same file, bottom-up passes and complexity.
\sourcenote{Edited from \emph{Small Proofs from Congruence Closure}, \S4.}
Our \textsc{TreeOpt} algorithm computes these costs from the input equalities upward.
An edge recording an input equality has cost one;
  an edge justified by congruence initially has infinite cost.
Each pass uses shortest paths between children to improve the costs
  of the corresponding congruence edges.
Repeated passes account for increasingly nested uses of congruence.
When no cost improves, a shortest path between the requested terms,
  together with its subproofs, gives a proof of minimum tree size.

% Source: same file, final paragraph; c = |C| and m = |E|.
\sourcenote{Edited from \emph{Small Proofs from Congruence Closure}, \S4;
  dependence on the congruence count emphasized.}
Let $c$ be the number of congruence edges and $m$ the total number of edges
  in the proof graph, including congruence edges.
\textsc{TreeOpt} needs at most $c$ passes,
  each updating the costs of at most $c$ congruence edges.
Its running time is $O(c^2m)$~\cite{eggproof}.
This parameterized bound exposes the cost of congruence reasoning:
  more congruence justifications require both more subproof searches per pass
  and potentially more passes.
\textsc{TreeOpt} establishes a polynomial-time solution for the relaxed tree-size objective,
  but these repeated searches make it too expensive for practical proof generation.
Its role is to provide a theoretical result, a baseline for evaluation,
  and a starting point for the greedy algorithm.

\subsubsection{Greedy Proof Search}
\label{sec:greedy-proofs}

% Source: ../egg-papers/2022-explanations-fmcad/05-flattened-greedy.tex,
% Greedy Optimization and Estimating Tree Sizes.
\sourcenote{Edited from \emph{Small Proofs from Congruence Closure}, \S5.}
Practical greedy proof search replaces optimal subproof costs with estimates
  from the trees maintained by traditional proof generation.
These trees already provide valid proofs, whose tree sizes bound the optimal costs.
Using the estimates as edge weights, the algorithm first finds a short path
  between the requested terms, then recursively improves the subproofs
  required by its congruence edges.
It limits this recursive optimization to a fixed number of subproof searches
  and uses the original tree proofs for the remaining subproofs.

% Source: same file, final complexity analysis and implementation limits.
\sourcenote{Edited from \emph{Small Proofs from Congruence Closure}, \S5.}
The implementation controls the expensive searches identified by \textsc{TreeOpt}'s analysis.
Let $n$ count the terms and input equalities in the instance.
For a fixed budget of recursive searches, computing the estimates and generating a proof
  takes $O(n\log n+c)$ time.
Thus greedy proof search is also parameterized by $c$, the number of congruence edges.
Our implementation permits ten recursive subproof searches
  and retains at most $10n$ congruence edges.
This bounds $c$ by a linear function of $n$ and gives an $O(n\log n)$ running time,
  matching the asymptotic bound for congruence closure.
The greedy algorithm guarantees a valid proof, but does not guarantee
  minimum tree size or minimum DAG size.

\sourcenote{New scope note for this report.}
For space, we summarize the algorithms' main ideas;
  their complete descriptions and analyses appear in the paper~\cite{eggproof}.

\subsection{Evaluation and Applications}
\label{sec:proofs-evaluation}

% Source: Small Proofs, 06-eval.tex, Z3 comparison and Detailed Analysis.
\sourcenote{Edited from \emph{Small Proofs from Congruence Closure}, \S6.}
We evaluated \texttt{egg}'s proof generation on 3,760 input/output pairs
  recorded from Herbie's equality saturation runs~\cite{eggproof}.
We asked Z3 to prove the same equalities from the same rewrite rules.
Z3 timed out after two minutes on 5\% of the queries;
  the proof-size comparison uses the remaining 3,571 pairs.
Greedy proofs were 72.8\% as large as Z3's on average, measured by DAG size.
They were also smaller than proofs obtained by removing unnecessary equalities
  from an existing \texttt{egg} certificate, a technique called \emph{proof reduction}.
This comparison shows that the improvement does not depend solely
  on differences between the \texttt{egg} and Z3 implementations.

% Source: ../egg-papers/2022-explanations-fmcad/06-eval.tex, Detailed Analysis and Table I.
\sourcenote{Edited from \emph{Small Proofs from Congruence Closure}, \S6.}
We use \textsc{TreeOpt} as an exact baseline for the relaxed tree-size objective.
Greedy proof search matched this optimum in 81.1\% of cases.
On the 3,571-query comparison, \texttt{egg} with greedy proof generation
  took 39.33\,ms per query on average, compared with 1,008.60\,ms with \textsc{TreeOpt}.
The greedy algorithm thus retains much of the benefit at a lower cost.
These results do not establish how close its proofs are to minimum DAG size:
  computing that optimum was too expensive to provide a baseline.

% Source: same evaluation, hardware case study and Table II.
\sourcenote{Edited from \emph{Small Proofs from Congruence Closure}, \S6, hardware case study.}
To measure the benefit for validation, we applied the algorithm to six Intel circuit benchmarks.
The validation flow checks each proof step with the industrial HECTOR equivalence checker;
  smaller proofs reduce the number of intermediate circuits it must compare.
DAG size fell by 29\% on average.
For a constant-multiplication design, the complete verification flow fell from 4.7 to 2.3 hours.
This demonstrates that smaller certificates can make an optimizer's results substantially cheaper to use.

% New connection between existing proof minimization and the proposed library.
\sourcenote{New connection to the proposed compiler.}
The proposed library in Section~\ref{sec:verified-compilation} extends this use of certificates
  to proofs containing analysis and renaming steps.
Its evaluation must determine whether minimizing those richer proofs
  similarly reduces runtime validation costs.

\subsection{Related Work}
\label{sec:proofs-related-work}

% Source: Small Proofs, 02-background.tex, Shrinking Congruence Proofs.
\sourcenote{Edited from \emph{Small Proofs from Congruence Closure}, \S2.}
Proof reduction removes equalities from an existing certificate
  and checks whether the remainder still proves the result~\cite{pp-congr}.
Trying each removal from a proof with $k$ input equalities takes $O(k^2\log k)$ time.
It cannot recover a shorter explanation using equalities already discarded.
Our algorithms consider alternatives during proof generation, before committing to a certificate.
